The problem file defines, for fixed $n,t$ and $0\le \alpha<1/2$, a parameter $F^{(t)}(n,\alpha)$ related to $2$-colorings of the complete $t$-uniform hypergraph on $n$ vertices.

\paragraph{Ambiguity/typo check.}
As written in the file:
\begin{quote}
``Let $F^{(t)}(n,\alpha)$ be the \emph{largest} $m$ such that we can $2$-colour the edges ... such that if $X\subseteq[n]$ with $|X|\ge m$ then there are at least $\alpha\binom{|X|}{t}$ many $t$-subsets of $X$ of each colour.''
\end{quote}
This makes $F^{(t)}(n,\alpha)$ trivial (equal to $n$) for every $\alpha\le 1/2$ (Lemma 161.1 below), contradicting the subsequent discussion in the same problem.

\medskip
\noindent\textbf{Minimal correction consistent with the surrounding text.}
Replace ``largest'' by ``smallest.'' Namely, define $F^{(t)}(n,\alpha)$ as the \emph{smallest} $m$ for which there exists a $2$-coloring of $\binom{[n]}{t}$ such that every subset $X\subseteq[n]$ with $|X|\ge m$ contains at least $\alpha\binom{|X|}{t}$ edges of each color.
This corrected parameter is monotone nondecreasing in $\alpha$ and can exhibit ``jumps.''
The question then becomes: as $\alpha$ increases from $0$ to $1/2$, is $F^{(t)}(n,\alpha)$ continuous in $\alpha$ (after rescaling) or does it have jumps (and how many)?
